\chapter{Εισαγωγή}

\section{Επιστημονική Περιοχή της Διπλωματικής}
Η ραγδαία εξέλιξη του Διαδικτύου των Πραγμάτων (\textlatin{Internet of Things - IoT}) και των υπολογιστικών υποδομών παρυφής (\textlatin{Edge Computing}) έχει οδηγήσει στην ανάδυση της τεχνολογίας των Ψηφιακών Διδύμων (\textlatin{Digital Twins}). Στη σύγχρονη βιβλιογραφία, ένα Ψηφιακό Δίδυμο ορίζεται ως η δυναμική, εικονική αναπαράσταση ενός φυσικού αντικειμένου, συστήματος ή διαδικασίας, η οποία τροφοδοτείται με δεδομένα πραγματικού χρόνου και επιτρέπει την προσομοίωση, την παρακολούθηση και, το σημαντικότερο, τον αμφίδρομο έλεγχο \cite{Grieves2017}, \cite{Han2022}.

Η μετάβαση από το απλό 'Ψηφιακό Μοντέλο' (\textlatin{Digital Model}) στην 'Ψηφιακή Σκιά' (\textlatin{Digital Shadow}) και, τελικά, στο πλήρες 'Ψηφιακό Δίδυμο', αποτελεί αντικείμενο εκτενούς έρευνας. Παράλληλα, σε ευρωπαϊκό επίπεδο, πρωτοβουλίες όπως το πρόγραμμα \textit{\textlatin{DIGITAL EUROPE}} επενδύουν στρατηγικά στη δημιουργία Τοπικών Ψηφιακών Διδύμων (\textlatin{Local Digital Twins}) για την ενίσχυση της βιωσιμότητας και της λειτουργικότητας των έξυπνων κοινοτήτων \cite{EuropeanCommission2023}. Η παρούσα διπλωματική εργασία εντάσσεται στον ευρύτερο ερευνητικό τομέα των Έξυπνων Κτιριακών Υποδομών (\textlatin{Smart Buildings}) και ειδικότερα στην εξειδίκευση των Έξυπνων Πανεπιστημιουπόλεων και Σχολείων (\textlatin{Smart Campuses}), μελετώντας την ενσωμάτωση του IoT στην εκπαιδευτική διαδικασία και τη λειτουργική διαχείριση (\textlatin{Facility Management}) \cite{Balla2023, Barros2024, Evangelou2022, Cui2023, EUParliament2024}.

\section{Σκοπός και Συνεισφορά της Διπλωματικής}
Σκοπός της παρούσας διπλωματικής εργασίας είναι ο σχεδιασμός, η υλοποίηση και η τεχνοοικονομική αξιολόγηση μιας ολοκληρωμένης πλατφόρμας Ψηφιακού Διδύμου, προσαρμοσμένης αποκλειστικά για σχολικά κτίρια. Ως πεδίο εφαρμογής (\textlatin{Case Study}) επιλέχθηκαν τα «Εκπαιδευτήρια Πλάτων» στην Κατερίνη.

Σε αντίθεση με συμβατικά συστήματα τηλεμετρίας, το προτεινόμενο σύστημα εγκαθιδρύει έναν κλειστό βρόχο αλληλεπίδρασης (\textlatin{Sense - Decide - Act}) μεταξύ του φυσικού και του ψηφιακού χώρου. Το σύστημα αναπτύχθηκε με χρήση ανοιχτού κώδικα (\textlatin{openHAB} ως \textlatin{semantic middleware}) και IoT εξοπλισμό χαμηλού κόστους (όπως μικροελεγκτές \textlatin{ESP32}, έξυπνους μετρητές \textlatin{Shelly 1PM} και ρελέ \textlatin{Sonoff}), που επικοινωνούν μέσω του πρωτοκόλλου \textlatin{MQTT}. Το σύστημα εξελίχθηκε σε ένα \textlatin{Enterprise Cloud SaaS} περιβάλλον (\textlatin{React 18}, \textlatin{Supabase}), το οποίο παρέχει την τελική εμπειρία χρήστη με \textlatin{responsive} πρόσβαση, \textlatin{3D} οπτικοποίηση και πολυεπίπεδο έλεγχο πρόσβασης.

Η κύρια συνεισφορά της εργασίας εστιάζεται σε τρεις άξονες:
\begin{itemize}
\item \textbf{Ποιότητα Εσωτερικού Περιβάλλοντος (\textlatin{IEQ}) και Μάθηση:} Η δυναμική παρακολούθηση της ακουστικής άνεσης στις αίθουσες διδασκαλίας (διατηρώντας τα επίπεδα κάτω από το όριο των 55~\textlatin{dBA}) με χρήση διακριτικής οπτικής ανατροφοδότησης (\textlatin{quiet-lamp}) προς τους μαθητές, προστατεύοντας την εκπαιδευτική διαδικασία \cite{WHO1999, Kephalopoulos2015}.
\item \textbf{Ενεργειακή Αποδοτικότητα:} Η βελτιστοποίηση της λειτουργίας ηλεκτροβόρων συσκευών, κυρίως στον χώρο του κυλικείου, που οδήγησε σε επαληθευμένη μείωση της ενεργειακής κατανάλωσης κατά 19,7\%.
\item \textbf{Λειτουργική Διαχείριση (\textlatin{Facility Management}):} Η ανάπτυξη ενός κεντρικού συστήματος τρισδιάστατης οπτικοποίησης (\textlatin{React 3D Dashboard}) και παρακολούθησης της ασφάλειας (π.χ. κατάσταση θυρών ψυγείων, έλεγχος πρόσβασης πυλών), επιτρέποντας την προγνωστική συντήρηση.
\end{itemize}

Η βασική ερευνητική υπόθεση (\textlatin{research hypothesis}) της παρούσας μελέτης είναι ότι η μετάβαση από ένα παθητικό σύστημα παρακολούθησης σε ένα ενεργό Ψηφιακό Δίδυμο (\textlatin{Digital Twin}) μπορεί να επιφέρει άμεση και μετρήσιμη βελτίωση στις λειτουργικές παραμέτρους μιας σχολικής υποδομής. Ειδικότερα, εξετάζεται (α) εάν ο αμφίδρομος έλεγχος μπορεί να διατηρήσει τα επίπεδα θορύβου κάτω από το όριο των 55~\textlatin{dBA} προστατεύοντας τη μαθησιακή διαδικασία, (β) εάν επιτυγχάνεται σημαντική μείωση της μέσης ημερήσιας κατανάλωσης ενέργειας (\textlatin{kWh}/ημέρα), και (γ) εάν η τεχνοοικονομική ανάλυση αποδεικνύει υψηλή απόδοση επένδυσης (\textlatin{ROI}) και ταχεία απόσβεση, επιβεβαιώνοντας την εμπορική βιωσιμότητα της πλατφόρμας.

Η αναπτυχθείσα πλατφόρμα είναι διαθέσιμη σε λειτουργία πραγματικού χρόνου στη διεύθυνση \href{https://digitaltwin.pi-tech.gr}{\textlatin{digitaltwin.pi-tech.gr}}, παρέχοντας δημόσια πρόσβαση σε επιλεγμένες προβολές της αρχιτεκτονικής και των ευρημάτων της έρευνας, καθώς και πλήρη λειτουργική πρόσβαση μέσω διαπιστευτηρίων για εξουσιοδοτημένους χρήστες. Μια επισκόπηση της κεντρικής διεπαφής της πλατφόρμας παρουσιάζεται στο Σχήμα~\ref{fig:dashboard_overview}.

\begin{figure}[htbp]
\centering
\includegraphics[width=\textwidth]{fig_1_1}
\caption{Επισκόπηση της διεπαφής της προτεινόμενης πλατφόρμας Ψηφιακού Διδύμου}
\label{fig:dashboard_overview}
\end{figure}

\section{Διάρθρωση της Διπλωματικής}
Η παρούσα διπλωματική εργασία αναπτύσσεται σε οκτώ (8) κεφάλαια, τα οποία διαρθρώνονται ως εξής:

Στο \textbf{Κεφάλαιο 1} (παρόν κεφάλαιο) γίνεται η εισαγωγή στο αντικείμενο, καθορίζεται ο σκοπός και παρουσιάζεται η δομή της μελέτης.

Στο \textbf{Κεφάλαιο 2} αναλύεται το θεωρητικό υπόβαθρο και το συναφές έργο γύρω από τα Ψηφιακά Δίδυμα στην εκπαίδευση, την Ποιότητα Εσωτερικού Περιβάλλοντος και τα πρωτόκολλα επικοινωνίας \textlatin{IoT}.

Το \textbf{Κεφάλαιο 3} περιγράφει την Αρχιτεκτονική του Συστήματος βάσει του μοντέλου των 5 επιπέδων, παρουσιάζοντας τον σχεδιασμό του υλικού (\textlatin{Hardware}) και τον διαχωρισμό \textlatin{Digital Shadow} και \textlatin{Digital Twin}.

Το \textbf{Κεφάλαιο 4} εμβαθύνει στην Τεχνική Υλοποίηση του λογισμικού, αναλύοντας τις δύο φάσεις υλοποίησης: το τοπικό πρωτότυπο (\textlatin{openHAB} / \textlatin{InfluxDB} / \textlatin{Grafana}) και την \textlatin{Enterprise Cloud SaaS} πλατφόρμα (\textlatin{React 18}, \textlatin{Supabase}, \textlatin{3D WebGL}).

Στο \textbf{Κεφάλαιο 5} παρουσιάζεται η μεθοδολογία αξιολόγησης του συστήματος στο πραγματικό περιβάλλον του σχολείου, εξηγώντας τον τρόπο μέτρησης της ενέργειας, της ακουστικής και της τεχνοοικονομικής ανάλυσης.

Το \textbf{Κεφάλαιο 6} παραθέτει τα πρακτικά αποτελέσματα (μείωση κατανάλωσης, εξοικονόμηση, βελτίωση ηχητικού περιβάλλοντος) και την τεχνοοικονομική μελέτη (Υπολογισμός \textlatin{ROI}).

Το \textbf{Κεφάλαιο 7} αφορά τη συζήτηση των ευρημάτων, προχωρώντας σε κριτική αξιολόγηση και ερμηνεία των αποτελεσμάτων.

Τέλος, στο \textbf{Κεφάλαιο 8} συνοψίζονται τα τελικά συμπεράσματα της έρευνας, καταγράφονται οι περιορισμοί και προτείνονται κατευθύνσεις για μελλοντική επέκταση της πλατφόρμας.





